% Part: incompleteness
% Chapter: representability-in-q
% Section: beta-function

\documentclass[../../../include/open-logic-section]{subfiles}

\begin{document}

\olfileid{inc}{req}{bet}
\olsection{बीटा फलनावरील साहाय्यक विधान}


बेरीज, गुणाकार आणि $\Char{=}$ उपलब्ध असतील तर आदिम पुनरावर्तन
करता येते हे दाखवण्यासाठी क्रमिका हाताळणारी फलने विकसित करावी लागतील.
(घातांक फलनही उपलब्ध असते, तर हे काम सोपे झाले असते.) आदिम
पुनरावर्तन उपलब्ध असताना ``$n$ वी अविभाज्य संख्या'' यांसारख्या
गोष्टी परिभाषित करून तुलनेने सरळ सांकेतीकरण निवडता येत होते.
पण येथे आदिम पुनरावर्तन उपलब्ध नाही---उलट, लघुतमीकरण वापरून
आदिम पुनरावर्तन करता येते हेच दाखवायचे आहे---म्हणून अधिक
युक्तीने काम करावे लागेल.

\begin{lem}
\ollabel{lem:beta}
असे फलन $\beta(d,i)$ आहे की प्रत्येक क्रमिका $a_0$,
\dots,~$a_n$ साठी अशी संख्या~$d$ असते की प्रत्येक $i \le n$
साठी $\beta(d,i) = a_i$. शिवाय, $\beta$ मूलभूत फलनांपासून
केवळ संयोजन आणि नियमित लघुतमीकरण वापरून परिभाषित करता येते.
\end{lem}

$d$ हा क्रमिका $\tuple{a_0, \dots, a_n}$ चा संकेतांक आहे
आणि $\beta(d,i)$ तिचा $i$ वा घटक देतो, असे समजा.
(लक्षात घ्या, या ``सांकेतीकरणात'' आपल्याला परिचित असलेले
अविभाज्य संख्यांच्या घातांवर आधारित सांकेतीकरण \emph{वापरलेले नाही}!)
हे साहाय्यक विधान किमान आवश्यक तेवढेच सांगते: क्रमिका जोडता
येतात, त्यांना घटक शेवटी जोडता येतो किंवा संयोजन व नियमित
लघुतमीकरणाने परिभाषित फलने वापरून $a_0$, \dots,~$a_n$
यांपासून $d$ \emph{काढता} येतो, असेही ते म्हणत नाही. ते फक्त
एवढेच सांगते की प्रत्येक क्रमिकेसाठी काही संकेतांक असतो आणि
त्यातून घटक मिळवणारे एक ``विसांकेतीकरण'' फलन आहे.

$\beta$ हे चिन्हांकन गोडेल यांचे आहे. पुन्हा सांगायचे तर,
या साहाय्यक विधानाच्या पुराव्यातील कठीण भाग म्हणजे वरवर
मर्यादित वाटणारी साधने---केवळ संयोजन आणि लघुतमीकरण---वापरून
योग्य $\beta$ परिभाषित करणे; अर्थात बेरीज, गुणाकार आणि
$\Char{=}$ वापरता येतात. या विधानाचे अनेक पुरावे आहेत; पण
त्यांपैकी एक सुबोध मार्ग आजही गोडेल यांची मूळ पद्धत आहे.
त्या पद्धतीत सुनझीचे प्रमेय (परंपरेने ``चिनी शेष प्रमेय'')
हे संख्या-सैद्धान्तिक तथ्य वापरले जाते.

\begin{defn}
दोन नैसर्गिक संख्या $a$ आणि $b$ या \emph{सापेक्षतः अविभाज्य}
(सहमूळ) असतात, तर आणि तरच त्यांचा महत्तम साधारण
विभाजक~$1$ असतो; म्हणजे त्यांचा दुसरा कोणताही सामाईक विभाजक नसतो.
\end{defn}

\begin{defn}
धन भाजक $c>0$ साठी नैसर्गिक संख्या $a$ आणि $b$ या \emph{भाजक~$c$ समशेष},
$a \equiv b \mod c$, असतात, तर आणि तरच
$c \mid (a-b)$; म्हणजे $a$ आणि $b$ यांना $c$ ने भागल्यावर
शेष समान मिळतो.
%READERNOTE{SOL6-A070 — समशेषतेचा समान शेष हा अर्थ धन भाजकासाठी आहे. शून्य भाजकाचा totalized शेष शून्य घेतला तरी तो divisibility द्वारे दिलेल्या समशेषतेशी समतुल्य नसतो. म्हणून या व्याख्येत आणि सुनझीच्या प्रमेयात धन भाजक स्पष्ट केला. बीटा फलाच्या बांधणीत सर्व भाजक आधीपासून धन आहेत.}
\end{defn}

सुनझीचे प्रमेय असे आहे:
\begin{thm}
समजा $x_0$, \dots,~$x_n$ या धन नैसर्गिक संख्या परस्पर जोड्यांनी सापेक्षतः
अविभाज्य आहेत. $y_0$, \dots,~$y_n$ या कोणत्याही संख्या असू द्या.
तर अशी संख्या $z$ आहे की
\begin{align*}
z & \equiv y_0 \mod x_0 \\
z & \equiv y_1 \mod x_1 \\
& \vdots  \\
z & \equiv y_n \mod x_n.
\end{align*}
\end{thm}

सुनझीचे प्रमेय आपण पुढीलप्रमाणे वापरू: $x_0$, \dots,~$x_n$
या अनुक्रमे $y_0$, \dots,~$y_n$ पेक्षा मोठ्या असतील, तर
$z$ हा क्रमिका $\tuple{y_0, \dots,y_n}$ चा संकेतांक घेता येतो.
$y_i$ परत मिळवण्यासाठी $z$ ला $x_i$ ने भागून शेष घ्यायचा.
या सांकेतीकरणासाठी $x_0$, \dots,~$x_n$ यांची योग्य मूल्ये
शोधावी लागतील.

यासाठी पुढील दोन निरीक्षणे उपयोगी पडतील. $y_0$, \dots,~$y_n$
दिली असता, पुढीलप्रमाणे घ्या:
\begin{align*}
j &= \max(n, y_0 + 1, \dots, y_n + 1), \\
m &= \lcm(1,\dots,j),
\end{align*}
दुसऱ्या ओळीत लघुतम साधारण विभाज्य घेतला आहे.
आणि पुढीलप्रमाणे घ्या:
\begin{align*}
x_0 & = 1 + m \\
x_1 & = 1 + 2 \cdot m \\
x_2 & = 1 + 3 \cdot m \\
& \vdots  \\
x_n & = 1 + (n+1) \cdot m
\end{align*}
मग पुढील दोन गोष्टी खऱ्या आहेत:
\begin{enumerate}
\item\ollabel{rel-prime} $x_0,\dots,x_n$ या परस्पर जोड्यांनी
  सापेक्षतः अविभाज्य आहेत.
\item\ollabel{less} प्रत्येक $i$ साठी $y_i < x_i$.
\end{enumerate}
\olref{rel-prime} हे खरे का, ते पाहू. दोन भिन्न निर्देशांक
निवडा. $p$ ही अविभाज्य संख्या
असेल आणि $p \mid x_i$ व $p \mid x_k$ असतील, तर
$p \mid 1 + (i+1) m$ आणि $p \mid 1 + (k+1) m$.
मग $p$ त्यांच्या फरकालाही विभाजित करते:
\[
(1 + (i+1)m) - (1+ (k+1)m) = (i-k) m.
\]
$p$ ही $1 + (i+1)m$ ला विभाजित करते, म्हणून ती $m$ ला
विभाजित करू शकत नाही (अन्यथा पहिल्या भागाकारात शेष~$1$
उरेल). त्यामुळे $p$ ही $(i-k)m$ ला विभाजित करत असल्याने
$p$ ही $i-k$ लाही विभाजित करते. पण $\left|i-k\right|$ जास्तीत
जास्त~$n$ आहे, आणि $j \geq n$ असे घेतले आहे; म्हणून
पुन्हा $p \mid m$, हा विरोधाभास. म्हणून $x_i$ आणि
$x_k$ या दोन्हींना विभाजित करणारी अविभाज्य संख्या नाही.
\olref{less} ही अट सोपी आहे: $y_i < j \leq m < x_i$.

आता $\beta$ फलनावरील साहाय्यक विधान सिद्ध करू. लक्षात ठेवा,
$0$, उत्तरवर्ती फलन, बेरीज, गुणाकार, $\Char{=}$, प्रक्षेपणे,
आणि नियमित फलनांना लघुतमीकरण लागू करून व संयोजनाने
यांपासून परिभाषित केलेली कोणतीही फलने वापरता येतात.
एखाद्या संबंधाचे जातिबोधक फल असे परिभाषित करता येत असेल,
तर तो संबंधही वापरता येतो. आधीप्रमाणे, हे संबंध बूलीय
संयोग आणि परिबद्ध संख्यकीकरण यांखाली बंद आहेत हे दाखवता येते;
उदाहरणार्थ:
\begin{align*}
\fn{not}(x) & \defis \Char{=}(x,0)\\
\bmin{x \leq z}{R(x,y)} & \defis \umin{x}{(R(x,y) \lor x = z)}\\
\bexists{x \leq z}{R(x,y)} & \defiff R(\bmin{x \leq z}{R(x,y)}, y)
\end{align*}
मग पुढील सर्व गोष्टी आदिम पुनरावर्तनाशिवायही परिभाषित
करता येतात, हे दाखवता येते:
\begin{enumerate}
\item जोडीकरण फलन, $J(x,y) = \frac{1}{2}[(x+y)(x+y+1)] + x$;
% येथे आणखी स्पष्टीकरण देता येईल; मांडणी थोडी गुंतागुंतीची आहे.
\item प्रक्षेपण फलने
\begin{align*}
K(z) & = \bmin{x \leq z}{\bexists{y \leq z}{z = J(x,y)}},\\
L(z) & = \bmin{y \leq z}{\bexists{x \leq z}{z = J(x,y)}};
\end{align*}
\item $x < y$ हा लघुत्वाचा संबंध;
\item $x \mid y$ हा विभाज्यतेचा संबंध;
% \item $x \tsub y$
% \item $\fn{Prime}(x)$
% \item $p$ अविभाज्य आहे असे धरल्यास, ``$x$ हा $p$ चा घात आहे'' हा संबंध:
% \[
% \bforall{y \leq x}{(y \mid x \lif y = 1 \lor y = x)}.
% \]
\item $y$ ला $x$ ने भागल्यावर मिळणारा शेष देणारे
  $\fn{rem}(x,y)$ हे फलन.
\end{enumerate}
शून्य भाजकासाठी या शेष फलाचे मूल्य येथे दिलेले नाही; आधीच्या
शेष फलाच्या परिपाठाप्रमाणे ते शून्य घेता येते. पुढील
वापरातील भाजक मात्र नेहमी धन असतो.
आता पुढील व्याख्या करा:
\begin{align*}
\beta^*(d_0,d_1,i) & = \fn{rem}(1+(i+1) d_1,d_0) \text{ आणि}\\
\beta(d,i) & = \beta^*(K(d),L(d),i).
\end{align*}
हेच आपल्याला हवे असलेले फलन आहे. वर दिलेल्या
$a_0,\dots,a_n$ साठी
\[
j = \max(n,a_0+1,\dots,a_n+1),
\]
असे घ्या, आणि $d_1 = \lcm(1,\dots,j)$ असे घ्या.
वरील \olref{rel-prime} नुसार $1+d_1$, $1+2 d_1$,
\dots, $1+(n+1) d_1$ या सापेक्षतः अविभाज्य आहेत,
आणि \olref{less} नुसार त्या अनुक्रमे
$a_0,\dots,a_n$ पेक्षा मोठ्या आहेत. सुनझीच्या प्रमेयानुसार
प्रत्येक~$i$ साठी पुढील अट पूर्ण करणारे $d_0$ चे मूल्य आहे:
\[
d_0 \equiv a_i \mod (1+(i+1)d_1)
\]
म्हणून ($d_1$ हे~$a_i$ पेक्षा मोठे असल्यामुळे)
\[
a_i = \fn{rem}(1+(i+1)d_1,d_0).
\]
$d = J(d_0,d_1)$ असे घ्या. मग प्रत्येक $i \le n$ साठी
\begin{align*}
\beta(d,i) & =  \beta^*(d_0,d_1,i) \\
& =  \fn{rem}(1+(i+1) d_1,d_0) \\
& =  a_i
\end{align*}
हेच अपेक्षित होते. याने $\beta$-फलनावरील साहाय्यक विधानाचा
पुरावा पूर्ण होतो.

\begin{prob}
  $x < y$, $x \mid y$ हे संबंध आणि $\fn{rem}(x,y)$ हे
  फलन आदिम पुनरावर्तनाशिवाय परिभाषित करता येते, हे दाखवा.
  $0$, उत्तरवर्ती फलन, बेरीज, गुणाकार, $\Char{=}$,
  प्रक्षेपणे, तसेच परिबद्ध लघुतमीकरण आणि परिबद्ध संख्यकीकरण
  वापरू शकता.
\end{prob}

\end{document}
